% p2_determinant.tex -- reverse-CS / rank-1 decomposition of det K_2
% (round 436, PRIME.RDAGGER.P2_DETERMINANT.01).
% Builds with pdflatex.
% Companion machine check: verify_p2_determinant.py.
% Discovery: experiments/tfpt-discovery/p2_determinant_probe.py.
% NO RH CLAIM.  Research documentation, not a theorem of RH.
\documentclass[11pt]{article}

\usepackage[a4paper,margin=2.7cm]{geometry}
\usepackage{amsmath,amssymb,amsthm}
\usepackage{microtype}
\usepackage{booktabs}
\usepackage{xcolor}
\usepackage[colorlinks=true,linkcolor=blue!50!black,urlcolor=blue!50!black]{hyperref}

\newtheorem{lemma}{Lemma}
\newtheorem{prop}{Proposition}
\theoremstyle{remark}
\newtheorem{remark}{Remark}

\newcommand{\half}{\tfrac12}
\newcommand{\R}{\mathbb{R}}
\newcommand{\lam}{\lambda_{-}}
\newcommand{\psim}{\psi_{-}}
\newcommand{\Rp}{R_{+}}
\newcommand{\Kp}{K_{2,+}}
\newcommand{\indm}{\operatorname{ind}_{-}}

\title{\bfseries The P2 determinant as reverse Cauchy--Schwarz\\[2mm]
\large Rank-1 identity, contribution anatomy, and the coarse
bound that dies at one window}
\author{Stefan Hamann}
\date{August 30, 2026}

\begin{document}
\maketitle

\begin{abstract}\noindent
On a real Hermitian invertible matrix $A_{0}$ with exactly one
negative eigenvalue $\lam<0$ and a two-column factor $U=(u_{1},u_{2})$,
the Haynsworth matrix $K_{2}=I_{2}+U^{T}A_{0}^{-1}U$ admits the
exact reverse-Cauchy identity
\[
\det K_{2}
=\det(\Kp)-Q,
\qquad
Q=p^{T}\operatorname{adj}(\Kp)\,p,
\]
where $\Kp=I+\Rp$ is the positive remainder (the $J^{-1}=I_{2}$
summand plus the positive-rest Gram) and
$p=U^{T}\psim/\sqrt{|\lam|}$.  Equivalently
$Q=\|p\|^{2}+p^{T}\operatorname{adj}(\Rp)p$.
The sign $\det K_{2}<0$ is $Q>\det(\Kp)$, i.e.\ the rank-1
negative-mode update dominates the remainder Gram.
A direction-free sufficient bound $\|p\|^{2}>\lambda_{\max}(\Kp)$
holds on $44/45$ of the nneg-$1$ family and is killed at the
same window that killed the r375 Cauchy test (kz\,46).
The contribution anatomy is mixed (no single term carries
$-\det K_{2}$).  The $C_{\min}$ mode of the r407 intertwiner
\emph{is} $\psim$.  Dropping $\psim$ restores $\det>0$;
vacuous windows ($n_{-}=0$) have no pole and $\det K_{2}>0$.
The r377 dictionary $\det K_{2}<0\iff h_{N}h_{N+1}<0$ is the
LDL reading of the same mixed $2\times2$.
What remains of (P2) is the dressed coupling
$p^{T}\Kp^{-1}p>1$ of the last two dual CD columns against
$\psi_{C_{\min}}$.  Nothing here is evidence for or against
the Riemann hypothesis.
\end{abstract}

\medskip\noindent
\textbf{Status.}\enspace
Lemmas~\ref{lem:id}--\ref{lem:suff} and~\ref{lem:cmin} are
proved (finite real algebra).
Lemma~\ref{lem:alt} is the r377 dictionary, re-read as the
sign of the post-cap LDL pair.
Proposition~\ref{prop:anatomy} is a finite machine check on the
$45$ resolvable nneg-$1$ windows, not a cofinal law.
The direction-free coarse bound is \emph{not} a theorem of the
family (kz\,46).
The original statement $\det K_{2}<0$ on the nneg-$1$ branch
is reduced to the dressed overlap $p^{T}\Kp^{-1}p>1$, with
source ingredient $\langle u_{\mathrm{CD}},\psi_{C_{\min}}\rangle$.
No RH claim.  No L$^{*}$ claim.  No R$^{\dagger}$ claim.

\section{The object}\label{sec:obj}

Let $A_{0}$ be real symmetric and invertible, of size $m\ge 2$,
with inertia $(m-1,1,0)$: one simple negative eigenvalue $\lam<0$
and a positive definite complement.  Write $\psim$ for a unit
eigenvector of $\lam$, $A_{+}^{-1}$ for the inverse of $A_{0}$ on
$\psim^{\perp}$ (zero on $\operatorname{span}\{\psim\}$), and
$U=(u_{1},u_{2})\in\R^{m\times 2}$.  Set
\[
K_{2}=I_{2}+U^{T}A_{0}^{-1}U.
\]
In the dual CD construction of rounds 367--371 one has
$A_{0}=R_{N-3}-\half I$ on the node set $Y$ and $U$ the last two
dual Christoffel--Darboux columns.  The Haynsworth identity
(sorry-free in \texttt{RH/Haynsworth.lean}) converts
$\indm(A_{0})=1$ together with $\det K_{2}<0$ into
$A_{0}+UU^{T}\succ 0$.  The mixed $3\times 3$ of round 369 is
$J=\operatorname{diag}(1,1,1/\mathrm{den})$; the present $2\times 2$
is the case $J=I_{2}$, so the $J^{-1}$ contribution is exactly $I_{2}$.

\section{Exact reverse-CS identity}\label{sec:id}

Write $w=U^{T}\psim$ and
$p=w/\sqrt{|\lam|}$, so $p\in\R^{2}$.  The spectral calculus is
$A_{0}^{-1}=\lam^{-1}\,\psim\psim^{T}+A_{+}^{-1}$ with $\lam<0$, hence
\[
U^{T}A_{0}^{-1}U
=-pp^{T}+\Rp,
\qquad
\Rp=U^{T}A_{+}^{-1}U\succeq 0.
\]
Therefore
\begin{equation}\label{eq:split}
K_{2}
=\Kp-pp^{T},
\qquad
\Kp=I_{2}+\Rp.
\end{equation}
The $I_{2}$ is the $J^{-1}$ summand; $\Rp$ is the positive-rest Gram.

\begin{lemma}[Rank-1 determinant]\label{lem:id}
\[
\det K_{2}
=\det(\Kp)-Q,
\qquad
Q=p^{T}\operatorname{adj}(\Kp)\,p
=a_{+}p_{2}^{2}-2b_{+}p_{1}p_{2}+d_{+}p_{1}^{2},
\]
where $\Kp=\bigl(\begin{smallmatrix}a_{+}&b_{+}\\ b_{+}&d_{+}\end{smallmatrix}\bigr)$.
Equivalently $Q=\det(\Kp)\,p^{T}\Kp^{-1}p$, and
\begin{equation}\label{eq:Jsplit}
Q
=\|p\|^{2}+p^{T}\operatorname{adj}(\Rp)\,p.
\end{equation}
On the rational $4$-node model
$A_{0}=\operatorname{diag}(-1,1,2,3)$,
$U=[(2,0,0,0)^{T},(1,1,0,0)^{T}]$ one has $p=(2,1)$,
$\Kp=\operatorname{diag}(1,2)$, $\det(\Kp)=2$, $Q=9$,
$\det K_{2}=-7$, and~\eqref{eq:Jsplit} reads $5+4$.
\end{lemma}

\begin{proof}
The matrix-determinant lemma on~\eqref{eq:split} is
$\det(\Kp-pp^{T})=\det(\Kp)\,(1-p^{T}\Kp^{-1}p)$.
The $2\times 2$ identity
$\det(\Kp)\,p^{T}\Kp^{-1}p=p^{T}\operatorname{adj}(\Kp)p$
is the first display.
For the split, $\operatorname{adj}(I+\Rp)=I+\operatorname{adj}(\Rp)$
as $2\times 2$ matrices, which is~\eqref{eq:Jsplit}.
The rational model is direct arithmetic (companion script).
\end{proof}

\begin{lemma}[Sign carrier]\label{lem:sign}
Assume $\Kp\succ 0$ (true once $\Rp\succeq 0$).  Then
\[
\det K_{2}<0
\quad\Longleftrightarrow\quad
Q>\det(\Kp)
\quad\Longleftrightarrow\quad
p^{T}\Kp^{-1}p>1.
\]
The first factor $\det(\Kp)=\det(I+\Rp)\ge 1$ cannot change sign
(r375, SATZ).
\end{lemma}

\begin{proof}
Lemma~\ref{lem:id} plus $\det(\Kp)>0$.
\end{proof}

\begin{lemma}[Direction-free sufficient bound]\label{lem:suff}
If $\|p\|^{2}>\lambda_{\max}(\Kp)$, then $\det K_{2}<0$.
\end{lemma}

\begin{proof}
$Q=p^{T}\operatorname{adj}(\Kp)p\ge\lambda_{\min}(\operatorname{adj}(\Kp))\|p\|^{2}
=\det(\Kp)\|p\|^{2}/\lambda_{\max}(\Kp)$.
The hypothesis forces $Q>\det(\Kp)$.
\end{proof}

The bound is sufficient, not necessary.  It is the reverse-CS
form of the r375 Cauchy test
$\|w\|^{2}>|\lam|\bigl(1+\lambda_{\max}(\Rp)\bigr)$.

\section{The $C_{\min}$ mode}\label{sec:cmin}

Round 407 gives the intertwiner $R=C(I+C)^{-1}$, hence
$A_{0}=\half(C-I)(C+I)^{-1}$ and
$A_{0}^{-1}=2(C+I)(C-I)^{-1}$.  The two matrices are rational
functions of each other, so they share eigenspaces.

\begin{lemma}[$C_{\min}$ is $\psim$]\label{lem:cmin}
The bottom eigenvector of $C$ is $\psim$.  In particular
$\indm(A_{0})=1$ if and only if $n_{C}:=\#\{\mu(C)<1\}=1$, and
the coupling of the last two CD columns against the defect is
\[
p_{i}
=\frac{\langle u_{i},\psi_{C_{\min}}\rangle}{\sqrt{|\lam|}}.
\]
\end{lemma}

\begin{proof}
Spectral calculus of the Cayley-type map
$\lambda=(\mu-1)/(2(\mu+1))$: $\mu<1$ iff $\lambda<0$, and the
eigenvectors coincide.  Live cosine $1-O(10^{-12})$ on every
tested window (companion).
\end{proof}

Both objects live at the cap: $U$ is the pair of dual OP columns
of degrees $N-3,N-2$, and $C_{\min}$ is the unique mode of the
rank-$(N-3)$ dual Gram that has crossed $1$.  The remaining
arithmetic is a lower bound on that overlap, dressed by $\Kp^{-1}$.

\paragraph{Gate protocol.}
Lemma~\ref{lem:id}--\ref{lem:suff} and~\ref{lem:cmin} are
unconditional finite algebra (no PNT, no RH).
The named remainder $p^{T}\Kp^{-1}p>1$ is a source statement
about dual CD columns versus the $C$-defect; it is not a
prime-number theorem and is not an RH claim.

\section{Anatomy, and the coarse kill}\label{sec:anatomy}

\begin{prop}[Branch check, not a law]\label{prop:anatomy}
On the $45$ resolvable nneg-$1$ windows of the r367 family
(core-$42$ $+$ r286-$15$ $+$ EXT3/4; EXT5/6 are the sign-census
class and were not rebuilt):
\begin{itemize}
\item
identity~\eqref{eq:split} holds at f64 grade
(max residual $1.6\cdot 10^{-10}$ on well-conditioned rows);
\item
$\det K_{2}<0$ on $45/45$, $-\det K_{2}\in[1.157,\,1126.389]$,
census floor $0.50$ on $45/45$;
\item
Rayleigh excess $p^{T}\Kp^{-1}p-1\in[0.275,\,3.295]$ on the
core (min at kz\,18); EXT excess $\ge 0.527$;
\item
the direction-free bound of Lemma~\ref{lem:suff} holds
$44/45$ and \emph{fails} at kz\,46
($\|p\|^{2}=2.852<\lambda_{\max}(\Kp)=3.656$, alignment
$0.984$, $\det K_{2}=-2.364$ still);
\item
$\|p\|^{2}>\lambda_{\min}(\Kp)$ holds $45/45$ with minimum gap
$0.316$ (kz\,46) --- a census inequality, \emph{not} sufficient
(the worst adj-direction still requires $\|p\|^{2}>\lambda_{\max}$);
\item
the dominant summand of $Q$ is mixed:
$T_{a}=a_{+}p_{2}^{2}$ on $25/45$,
$T_{d}=d_{+}p_{1}^{2}$ on $4/45$,
$T_{x}=-2b_{+}p_{1}p_{2}$ on $16/45$ (the cross term grows on
EXT).  No single-term bound covers the family;
\item
the $J^{-1}$ share $\|p\|^{2}/Q$ lies in $[0.22,0.59]$ on the
core: the $I_{2}$ contribution alone does not dominate
($\|p\|^{2}>\det(\Kp)$ on only $5/28$ of the core);
\item
$\cos(\psim,\psi_{C_{\min}})=1$ and $n_{C}=1$ on $45/45$.
\end{itemize}
Fourteen of the forty-two core windows are vacuous
($n_{-}=0$): there $\Kp=K_{2}$, $\det K_{2}>0$ on $14/14$,
and $n_{C}=0$.  Overload $n_{-}\ge 2$ is $0/74$.
\end{prop}

The proposition is a machine check.  The Reviewer's hope that a
coarse $O(1)$ bound on $\|p\|$ would close (P2) because
$-\det K_{2}$ is huge is answered honestly: the \emph{output}
reserve is $O(1)$ and flat, but the direction-free input bound
$\|p\|^{2}>\lambda_{\max}(\Kp)$ is not universal on the family.
The same window (kz\,46) already killed the r375 Cauchy form; the
alignment $0.984$ is why (P2) still holds there.

\section{Alternation}\label{sec:alt}

\begin{lemma}[Post-cap dictionary]\label{lem:alt}
On the nneg-$1$ branch, $\det K_{2}<0$ if and only if the
product of the two LDL pivots of $K_{2}$ is negative.
The r377 dictionary identifies those pivots, up to a positive
Jacobi gauge, with $(h_{N},h_{N+1})$ of the signed source, so
\[
\det K_{2}<0
\quad\Longleftrightarrow\quad
h_{N}\,h_{N+1}<0.
\]
This is not the r417 Woodbury-$\mathrm{sch}$ coordinate
(r417: $\mathrm{sch}$ is a different $1\times 1$).
\end{lemma}

The $\psim$-summand is the sign carrier: without it,
$K_{2}=\Kp\succ 0$, both LDL pivots are positive, and
$h_{N}h_{N+1}>0$ (vacuous $29/29$; explicit kz\,12:
$\det K_{2}=+12.463$, first negative pivot at $N+2$).
On the branch the rank-1 update makes $K_{2}$ indefinite
iff $Q>\det(\Kp)$, which is exactly the alternation.
Live: $45/45$ nneg-$1$ alternate; $29/29$ vacuous do not.

\section{Kills}\label{sec:kills}

\begin{center}
\begin{tabular}{llll}
\toprule
mutant & hypothesis & $\det K_{2}$ & note \\
\midrule
vacuous $n_{-}=0$ (14/42 core) & empty & $>0$ (14/14) & no pole, $K_{2}=\Kp$ \\
drop $\psim$ at w9 & $n_{-}=1$ kept, pole omitted & $+5.958$ & P2 dies \\
wrong $J=-I$ at w9 & wrong remainder & $+0.941$ & sign flips \\
dead $\chi_{3}$-15 & $n_{-}=1$ & $-122.91$ & P2 \emph{stays} \\
scramble w9 & $n_{-}=21$ & $-8.881$ & hypothesis empty \\
tiny overlap (toy) & $n_{-}=1$, $\|p\|$ small & $+1.0025$ & P1 $\not\Rightarrow$ P2 \\
\bottomrule
\end{tabular}
\end{center}

Dead character-twisted rungs on the nneg-$1$ branch keep
$\det K_{2}<0$ (CHI3-$15/19/39$; CHI4-w9).  Terminal death is a
$q^{\dagger}$ statement (r433), not a P2 failure.  The two
coordinates are kept distinct from r417 $\mathrm{sch}$.

\section{Named remainder}\label{sec:rem}

Lemmas~\ref{lem:id}--\ref{lem:sign} convert (P2) into
\begin{equation}\label{eq:red}
p^{T}\Kp^{-1}p>1
\tag{P2$^{-}$}
\end{equation}
with $p_{i}=\langle u_{i},\psi_{C_{\min}}\rangle/\sqrt{|\lam|}$.
This is the r375 dressed overlap in reverse-CS letters, now
with an identified source vector $\psi_{C_{\min}}$ and with a
documented kill of every coarser direction-free bound that
was tried.  A numerical $\varepsilon$ is not frozen: kz\,18 has
excess $0.275$, compatible with any $\varepsilon\le 1/4$, but
a fitted constant is not a source statement.

The last two dual CD columns, against the unique $C$-mode below
$1$, must have $\Kp$-dressed overlap strictly larger than one.
Expanding that overlap via the dual CD kernel on $Y$ is the
remaining finite identity.  No $\lambda_{\min}$ of a Schur block
of the full $R$, and no terminal margin, enters the target.

\section{Claim boundary}

Finite real algebra on a Hermitian pair $(A_{0},U)$, plus a
named remainder on the dual CD construction.  Haynsworth remains
sorry-free and does not assert (P1) or (P2) on windows.
L$^{*}$, R$^{\dagger}\succeq\half I$, and the Riemann hypothesis
are not addressed in either direction.

\end{document}
